% Autonomous typing of the terminal dodecaphase chart.
\subsection{Signed dodecaphase chart of the full TPK state}
\label{subsec:registro-evaluacion-incidencial}

The dodecaphase chart and a possible factorization through the raw
108-step projection (CT108) are objects of different types. Let
\(X\in\mathsf{TPKFullState}\) be a complete enriched state, and let
\(\rho_{108}(X)\in\mathsf{CT108RawProjection}\) be the object obtained by
forgetting orientation, sheet, carry, boundary signature, memory, and the
dodecaphase ledger. In the canonical coordinates of the full state, write
\[
 X=(X_{\mathrm{APP}},X_{\mathrm{TRIT}},X_{\mathrm{route}},X_{\mathrm{mem}},
       C_{\mathrm{term}},U_{12}^{\mathrm{sgn}},
       X_{\mathrm{boundary}},\ldots).
\]
The active reader is, by definition, the projection of that already generated
and preserved coordinate:
\[
 \mathcal S_{12}:=\operatorname{pr}_{U_{12}^{\mathrm{sgn}}}:
 \mathsf{TPKFullState}\longrightarrow\mathsf{SignedDodecaphaseState}.
\]
On the compatible sublattice of the full state, this projection coincides
with the composite reader
\begin{equation}
 \mathcal S_{12}\equiv R_{\mathrm{sgn}}
 :=\Pi_H\circ\mathcal E_{108}^{90,120}\circ\mathcal R_{12}.
 \label{eq:registro-identidad-s12-rsgn}
\end{equation}
Therefore, \(\mathcal S_{12}\) reads a preserved coordinate of the same state;
it receives neither a conventional constant nor a metrological table. Its
domain is not replaced by \(\mathsf{CT108RawProjection}\), since that
replacement would additionally require a factorization
\(\mathcal S_{12}=\Xi\circ\rho_{108}\). Such a factorization is neither
presupposed nor needed here; it would require proof that the signed reading is
constant on the fibers of \(\rho_{108}\). The internal theorem operates
directly on the full TPK state.

To fix the content of a possible incidence realization, let \(\mathscr L\) be
a ledger of visits and oriented traces belonging to \(X\). Each visit
preserves the route, APP position, sheet, instant, multiplicity, and boundary
incidence. The arithmetic evaluation produces the integer and its
decomposition
\[
 n_v=r_v+9q_v,\qquad 1\leq r_v\leq9.
\]
The integer \(n_v\) is calculated from the sum or product sheet; its remainder
is not received as an independent value. The distinction between corona
incidence and interior incidence remains explicit for visits of remainder
\(9\).

In the temporal window
\(E_m=\{9(m-1)+1,\ldots,9m\}\), the domains are distinguished
\[
 \mathscr V_m=\{v:t(v)\in E_m\},\qquad
 \mathscr T_m=\{\theta:\operatorname{origin}(\theta)\in E_m\}.
\]
The presentation by counts considers
\begin{align}
 A_m&=\sum_{v\in\mathscr V_m}w(v)
             \mathbf1_{\{1,3,5,7\}}(r(v)),\\
 C_m&=-\sum_{\theta\in\mathscr T_m}u(\theta)\sigma(\theta)
      =\sum_{j\geq0}\bigl(\nu^-_{120,m}(j)-\nu^+_{120,m}(j)\bigr),\\
 V_m&=\sum_{v\in\mathscr V_m}w(v)
              \mathbf1_{\{9\}}(r(v))\epsilon_\partial(v).
\end{align}
Here \(w(v)\) is the incidence multiplicity, and
\(\epsilon_\partial(v)=1\) on the corona, \(-1\) in the interior.
The traces counted by \(\nu^\pm\) must be enumerated by the corresponding
ledger rule. The integral sum requires finite support or an alternative
construction that determines its meaning. The number of chronological
incidences and the reduced length \(120(1+3j)\) have different types; their
relation requires the reader's unit of length.

Define \([x]_{10}\in\{0,\ldots,9\}\) as the Euclidean representative and
preserve the quotients
\[
 x=[x]_{10}+10\lfloor x/10\rfloor.
\]
The incidence block is
\begin{equation}
 z_m=100[A_m]_{10}+10[C_m]_{10}+[V_m]_{10}.
 \label{eq:registro-bloque-incidencial}
\end{equation}
By composition one obtains the mapping
\begin{equation}
 \mathcal E_{\rm cnt}(\mathscr L)=
 \left((z_m-z_{m+3})_{m=1}^{12},
       (z_m-z_{m+4})_{m=1}^{12},\sum_{m=1}^{12}z_m\right).
\label{eq:registro-composicion-incidencial}
\end{equation}
The indices are read cyclically modulo \(12\). In this subsection,
\((D_jz)_m=z_m-z_{m+j}\), for \(j=3,4\), and
\(Q(z)=\sum_{m=1}^{12}z_m\); therefore, the preceding composition is
\((D_3z,D_4z,Q(z))\). This convention is preserved in
\S\ref{subsec:k-transversal}, where the complete integral reconstruction is
proved. This formula produces its coordinates from the counts. It uses
neither \(K\), \(U\), \(\alpha\), nor the expected transverse coordinates as
arguments.

\subsubsection{Terminal representation and exact closure of its charts}

The canonical terminal state is
\[
 X_{\mathrm{term}}
 =\bigl(\mathcal U_{108}\circ\cdots\circ\mathcal U_1\bigr)
   (X_{\mathrm{seed}}),
 \qquad
 \mathcal U_t=\mathsf{Upd}_t\circ\mathsf{Tra}_t\circ\mathsf{Sel}_t,
\]
where each \(\mathcal U_t\) acts on the full TPK state and not on its raw
projection. Among the coordinates produced and preserved before the signed
representation is
\[
 C_{\mathrm{term}}=\operatorname{pr}_{C}(X_{\mathrm{term}})
 =(b^{90},b^{120},Q_{\mathrm{TPK}}),
\]
with
\begin{align*}
 b^{90}={}&(-495,-116,-684,108,601,-90,
              -173,-88,313,560,-397,461),\\
 b^{120}={}&(-425,-281,-481,671,-255,30,
               475,-543,680,251,6,-128),\\
 Q_{\mathrm{TPK}}={}&6263.
\end{align*}
No subsequent table is introduced here: \(C_{\mathrm{term}}\) is the
transverse coordinate that \(\mathcal E_{108}^{90,120}\) produces within the
evolution of the full state. On the compatible sublattice of TPK states, the
chart satisfies the operator identity
\begin{equation}
 \mathcal S_{12}=\operatorname{pr}_{U}
 =\Pi_H\circ\operatorname{pr}_{C}.
 \label{eq:registro-identidad-s12-pih}
\end{equation}
The identity is formulated on the full TPK state and does not presuppose the
descent of \(\mathcal S_{12}\) through \(\rho_{108}\).

The subsequent evaluation from the twenty-five expressed fields is displayed
without a gap. For \(1\le r\le3\), let
\[
 B_r=\sum_{j=0}^{3}b^{120}_{r+3j},\qquad
 d_{r,j}=b^{90}_{r+3j},
\]
and
\[
 s_1=\frac{Q_{\mathrm{TPK}}+2B_1+B_2}{3},\qquad
 s_2=s_1-B_1,\qquad s_3=s_2-B_2.
\]
Integral substitution of the preceding twenty-five fields gives
\[
 (B_1,B_2,B_3)=(972,-1073,101),\qquad
 (s_1,s_2,s_3)=(2378,1406,2479).
\]
Applying
\[
 U^{(r)}=(s_r,d_{r,1}-d_{r,3},d_{r,0}+d_{r,2},d_{r,0}-d_{r,2})
\]
one obtains, component by component,
\begin{align*}
 U^{(1)}&=(2378,-452,-668,-322),\\
 U^{(2)}&=(1406,998,-204,-28),\\
 U^{(3)}&=(2479,-551,-371,-997).
\end{align*}
Columnwise interleaving therefore evaluates the signed chart:
\begin{equation}
 U_{\mathrm{term}}:=\mathcal S_{12}(X_{\mathrm{term}})
 =(2378,1406,2479,-452,998,-551,-668,-204,-371,-322,-28,-997).
 \label{eq:registro-uterm-desde-estado-pleno}
\end{equation}

\begin{theorem}[Subsequent exact closure among \texorpdfstring{\(U\)}{U}, \texorpdfstring{\(K\)}{K}, and the channels]
\label{thm:registro-evaluacion-terminal-recuperada}
Once the signed chart \eqref{eq:registro-uterm-desde-estado-pleno} has been
expressed by projection, the orbital Hadamard transform expresses, at a
subsequent stage,
\[
 K=(234,543,140,729,659,824,621,058,914,794,146,601),
\]
and the transverse extractor \((D_3,D_4,Q)\) produces
\begin{align}
 b^{90}={}&(-495,-116,-684,108,601,-90,
             -173,-88,313,560,-397,461),\notag\\
 b^{120}={}&(-425,-281,-481,671,-255,30,
              475,-543,680,251,6,-128),\notag\\
 Q_{\mathrm{TPK}}={}&6263.
 \label{eq:registro-canales-terminales-producidos}
\end{align}
The composition is an exact bijection among the compatible sublattice of signed
states, the seal \(K\), and its transverse coordinate
\(C_{\mathrm{term}}=(b^{90},b^{120},Q_{\mathrm{TPK}})\).
In particular,
\begin{equation}
 C_{\mathrm{term}}=(D_3K,D_4K,Q(K)).
 \label{eq:registro-cterm-desde-xterm}
\end{equation}
\end{theorem}

\begin{proof}
Grouping \(U_{\mathrm{term}}\) into its three step-three orbits yields
\[
 \begin{aligned}
  u^{(0)}&=(2378,-452,-668,-322),&
  \tfrac14H_4u^{(0)}&=(234,729,621,794),\\
  u^{(1)}&=(1406,998,-204,-28),&
  \tfrac14H_4u^{(1)}&=(543,659,58,146),\\
  u^{(2)}&=(2479,-551,-371,-997),&
  \tfrac14H_4u^{(2)}&=(140,824,914,601).
 \end{aligned}
\]
The divisions are exact in \(\mathbb Z\).
Columnwise interleaving expresses \(K\). Finally, the twelve cyclic
step-three and step-four subtractions and the total sum give, component by
component, \eqref{eq:registro-canales-terminales-producidos}. Proposition
\ref{prop:k-transversal} proves the integral inverse, and
Theorem~\ref{thm:k-naturalidad-lector-firmado} proves naturality under
refinement. No operation receives \(\alpha\), CODATA, a metrological table,
or a conventional constant.
\end{proof}

\begin{remark}[Complete operator form and focal executable control]
The terminal arrow is not postulated from a table: it is the typed
composition
\[
 X_{\mathrm{term}}
 \xrightarrow{\mathcal R_{12}}\mathscr L(X_{\mathrm{term}})
 \xrightarrow{\mathcal E_{108}^{90,120}}C_{\mathrm{term}}
 \xrightarrow{\Pi_H}U_{\mathrm{term}}
 \xrightarrow{\mathcal H_{12}^{-1}}K .
\]
The preceding definitions traverse the twelve windows that partition the one
hundred eight steps and express each channel coordinate as a finite sum of its
visits and traces, with orientation, sheet, carry, boundary, and memory. This
is the mathematical realization of the production of
\(C_{\mathrm{term}}\); a row-by-row tabular expansion would be another
serialization of the same sum, not an additional premise of the theorem. The
focal executable exactly reevaluates the finite chart
\(C_{\mathrm{term}}\leftrightarrow U_{\mathrm{term}}\leftrightarrow K\)
as a supplementary reproducible control. The proof of the continuum
hypothesis remains formulated over APP--TRIT--TPK, the full enriched state,
its correlated representations, and the joint discrete structure.
\end{remark}

Neither \(\mathcal S_{12}\), nor \(U\), nor \(K\), nor \(\alpha\) are
external inputs or conventional selectors of the discrete construction of the
continuum. They are typed coordinates and representations of the same
APP--TRIT--TPK state. The coordinate \(\alpha_{\mathrm{HMT}}\) is a derived
and correlated output of that single generation; its representation order does
not break the continuity of the object or weaken its derivation.

The incidence formula \eqref{eq:registro-composicion-incidencial} describes an
additional realization when a complete ledger
\(\mathscr L(X_{\mathrm{term}})\) is unfolded. The preceding theorem does not
claim that the signed coordinate can be reconstructed after that ledger has
been forgotten through \(\rho_{108}\); nor does it require such a claim. This
separation avoids confusing the internal closure of the enriched state with
an inverse problem posed on a projection that has removed precisely the
orientation and memory information.

The mapping \(\Pi_H\), executed within the chart, produces from
\(C_{\mathrm{term}}\) the signed coordinate and its three step-three orbits:
\begin{align*}
 U_{\mathrm{term}}^{(1)}&=\Pi_H^{(1)}(C_{\mathrm{term}})
                         =(2378,-452,-668,-322),\\
 U_{\mathrm{term}}^{(2)}&=\Pi_H^{(2)}(C_{\mathrm{term}})
                         =(1406,998,-204,-28),\\
 U_{\mathrm{term}}^{(3)}&=\Pi_H^{(3)}(C_{\mathrm{term}})
                         =(2479,-551,-371,-997).
\end{align*}
Let \(\mathcal P_{\mathrm{orb}}:\mathbb Z^{12}\to(\mathbb Z^4)^3\) be the
orbital regrouping
\[
 \mathcal P_{\mathrm{orb}}(U)
 =\bigl((U_1,U_4,U_7,U_{10}),
        (U_2,U_5,U_8,U_{11}),
        (U_3,U_6,U_9,U_{12})\bigr),
\]
and let \(\operatorname{Int}_{3,4}=\mathcal P_{\mathrm{orb}}^{-1}\) be the
columnwise interleaving of three tetrads. We also define
\[
 \operatorname{Dig}_{10}^{(3)}(z)
 =\left(\left\lfloor\frac z{100}\right\rfloor,
         \left\lfloor\frac z{10}\right\rfloor\bmod10,
         z\bmod10\right),\qquad0\le z\le999.
\]
The terminal modular chart is the effective mapping on the compatible channel
state
\begin{equation}
 \mathcal G_{\mathrm{term}}(C)
 :=\left[(\operatorname{Dig}_{10}^{(3)})^{12}
 \circ\operatorname{Int}_{3,4}
 \circ\left(\tfrac14H_4\oplus\tfrac14H_4
                    \oplus\tfrac14H_4\right)
  \circ\mathcal P_{\mathrm{orb}}\right]\!\left(\Pi_H(C)\right).
 \label{eq:registro-generador-modular-terminal}
\end{equation}
Its domain is the subset of the integrality sublattice of compatible
coordinates in \(\mathbb Z^{25}\) on which \(\Pi_H\) is integral, the three
orbital images under \(H_4\) are divisible by four, and the twelve resulting
coordinates belong to \(\{0,\ldots,999\}\). It receives neither
\(U_{\mathrm{term}}\), \(K\), \(\alpha\), a decimal table, nor a metrological
value; the interleaving, the phase origin, and the orientation are part of the
type of the channel coordinate.

The verification coordinate is serialized in
\path{metadata/semilla-registro-terminal.json}. That file contains
\(C_{\mathrm{term}}\), the phase origin, and the orientation; it contains
neither \(U_{\mathrm{term}}\), \(K\), nor the twelve residual rows. The focal
program applies \(\mathcal G_{\mathrm{term}}\), expresses
\(U_{\mathrm{term}}\), \(K\), and the residues, and returns through
\((D_3,D_4,Q)\) to the same \(C_{\mathrm{term}}\). Starting there permits an
exact focal verification of the subsequent finite chart. It does not redefine
the domain of the signed chart
\eqref{eq:registro-uterm-desde-estado-pleno}, does not replace the preceding
proof, and does not attribute to the auxiliary the production of the
coordinate it receives.

\begin{proposition}[Modular calculation of the twelve terminal rows]
\label{prop:registro-generacion-modular-terminal}
The evaluation of \(\mathcal G_{\mathrm{term}}\) on the produced coordinate
\(C_{\mathrm{term}}\) first obtains \(U_{\mathrm{term}}\) and then the phase
tetrads
\[
 (234,729,621,794),\qquad(543,659,058,146),\qquad
 (140,824,914,601),
\]
and, after interleaving, produces exactly the twelve rows of the following
table. It is the algebraic verification subsequent to
Theorem~\ref{thm:registro-evaluacion-terminal-recuperada}; it does not replace
the representation of the signed chart of the full TPK state.
\end{proposition}

\begin{proof}
The formulas for \(\Pi_H\) give the three signed blocks written above. The
subsequent product by \(H_4/4\) gives, without rounding,
\begin{align*}
 \tfrac14H_4(2378,-452,-668,-322)^T&=(234,729,621,794)^T,\\
 \tfrac14H_4(1406,998,-204,-28)^T&=(543,659,58,146)^T,\\
 \tfrac14H_4(2479,-551,-371,-997)^T&=(140,824,914,601)^T.
\end{align*}
The twelve divisions are exact in \(\mathbb Z\). Columnwise interleaving gives
\[
 \begin{split}
 (&234,543,140,729,659,824,\\
  &621,058,914,794,146,601),
 \end{split}
\]
and \(\operatorname{Dig}_{10}^{(3)}\) uniquely provides its three decimal
residues. The auxiliary verifier included in the package executes these same
operations; its negative controls alter the channel state or its phase origin.
\end{proof}

Denote by
\[
 \mathscr F_{\mathrm{term}}^{(10)}
 :=\bigl((a_m,c_m,v_m)\bigr)_{m=1}^{12}
\]
the digit chart expressed by \(\operatorname{Dig}_{10}^{(3)}(K_m)\).
When a complete incidence realization is available,
\((a_m,\allowbreak c_m,\allowbreak v_m)=([A_m]_{10},\allowbreak[C_m]_{10},\allowbreak[V_m]_{10})\).

This equality is not
postulated in order to reconstruct the ledger after it has been forgotten.
Each row preserves the window index and is therefore not a bag of
interchangeable digits.

\begin{center}
\begin{tabular}{c@{\qquad}ccc@{\qquad}c}
\toprule
\(m\)&\(a_m\)&\(c_m\)&\(v_m\)&\(K_m\)\\
\midrule
1&2&3&4&234\\
2&5&4&3&543\\
3&1&4&0&140\\
4&7&2&9&729\\
5&6&5&9&659\\
6&8&2&4&824\\
7&6&2&1&621\\
8&0&5&8&058\\
9&9&1&4&914\\
10&7&9&4&794\\
11&1&4&6&146\\
12&6&0&1&601\\
\bottomrule
\end{tabular}
\end{center}
Consequently,
\begin{equation}
 z^{\mathrm{term}}
 =(234,543,140,729,659,824,621,058,914,794,146,601).
 \label{eq:registro-terminal-residuos}
\end{equation}

\begin{proposition}[Terminal return of the computed collection to the channels]
\label{prop:registro-terminal-coordenadas}
The decimal recomposition of \(\mathscr F_{\mathrm{term}}^{(10)}\), followed
by the transverse extractor, produces, without consulting \(\alpha\) or a
metrological table,
\begin{align*}
 D_3z^{\mathrm{term}}={}&(-495,-116,-684,108,601,-90,
 &\hspace{18mm}-173,-88,313,560,-397,461),\\
 D_4z^{\mathrm{term}}={}&(-425,-281,-481,671,-255,30,
 &\hspace{18mm}475,-543,680,251,6,-128),\\
 Q(z^{\mathrm{term}})={}&6263.
\end{align*}
These are exactly the twenty-five coordinates
\((b^{90},b^{120},Q_{\mathrm{TPK}})\) used by the harmonic reader.
\end{proposition}

\begin{proof}
The decimal formation of each row gives
\eqref{eq:registro-terminal-residuos}. Cyclically subtracting the entry
\(m+3\) and the entry \(m+4\) produces, component by component, the two
vectors displayed; the sum of the twelve entries is \(6263\). The calculation
is repeated in the package's autonomous verifier. The terminal collection is a
computed representation of the channel state. The transverse evaluation returns
exactly to the seed \(C_{\mathrm{term}}\); the harmonic reading then recovers
\(U_{\mathrm{term}}\), and the inverse Hadamard transform recovers the same
\(K\). This proves the agreement of both coordinatizations without using
alpha as a selector.
\end{proof}

The chart \(\mathscr F_{\mathrm{term}}^{(10)}\) is sufficient and necessary
for this transverse evaluation. The table is a decimal certificate of the
expressed output, not a visit-by-visit reedition of the oriented ledger. The
enriched state additionally preserves Euclidean quotients, visits, traces,
and provenance; those fields are not declared recoverable from the thirty-six
residues, nor are they necessary for the cardinal theorem of this article.

\begin{proposition}[Sufficient and necessary arithmetic information]
\label{prop:registro-36-residuos}
Let \(N,N'\in\mathbb Z^{12\times3}\), ordered by the counts
\((A,C,V)\). For the mapping defined by
\eqref{eq:registro-bloque-incidencial}--\eqref{eq:registro-composicion-incidencial},
\[
 \mathcal E_{\rm cnt}(N)=\mathcal E_{\rm cnt}(N')
 \quad\Longleftrightarrow\quad
 N-N'\in10\mathbb Z^{36}.
\]
In particular, the \(36\) sectorial residues exactly determine the output of
\(25\) coordinates. The preservation of the quotients and provenance
constitutes additional information in the record.
\end{proposition}

\begin{proof}
Equality modulo \(10\) produces the same blocks and, therefore, the same
transverse coordinates. Conversely, if the transverse coordinates coincide,
\(D_3(z-z')=D_4(z-z')=0\). The steps \(3\) and \(4\) generate the cycle of
length \(12\), so \(z-z'\) is constant. Equality of the total charges makes
that constant vanish. Finally, the decimal representation of an integer
between \(0\) and \(999\) has three unique digits. The three residues of each
window are thus recovered. The assertion is an algebraic identity for every
\(N,N'\), not an inference obtained from finite tests.
\end{proof}

\subsection{Conjugation of the counts and boundary terms}
\label{subsec:registro-conjugacion-incidencial}

Let \(\rho(m)=13-m\). Consider an incidence conjugation that reverses the
order of the windows, preserves the arithmetic class and the boundary class
of the visits, and interchanges the trace channels. Then
\[
 (A_m^\dagger,C_m^\dagger,V_m^\dagger)
 =(A_{\rho(m)},-C_{\rho(m)},V_{\rho(m)}).
\]
If \(c_m=[C_m]_{10}\), decimal reduction produces the identity
\begin{equation}
 z_m^\dagger=z_{\rho(m)}
       +100\mathbf1_{\{c_{\rho(m)}\ne0\}}-20c_{\rho(m)}.
 \label{eq:registro-conjugacion-decimal}
\end{equation}
Indeed, \([-C]_{10}=0\) when \([C]_{10}=0\), and
\([-C]_{10}=10-[C]_{10}\) otherwise. The formula preserves exactly this
difference. Negating the channel is not equivalent to negating the complete
decimal block.

Applying this identity to a concrete TPK involution requires verifying its
action on the incidences. For an occupancy reader evaluated before the
advance, the reversal of a route
\(\gamma=(x_0,\ldots,x_n)\) satisfies
\[
 \sum_{k=0}^{n-1}f(x_{n-k})
 =\sum_{k=0}^{n-1}f(x_k)+f(x_n)-f(x_0).
\]
The endpoint terms are incorporated before reduction modulo \(10\). They
vanish on closed routes; in open windows they may alter the residues. This
correction, already present in the bidirectional development of transport,
specifies what the linkage between the route and the counts must preserve.

\begin{proposition}[Necessity of the nonperiodic information in the record]
Suppose that the observable of a fixed collection of routes satisfies
\(o(t+54)=o(t)\), that multiplicity is preserved under that return, and that
the same local reader acts on the windows of \(9\) instants. Then its counts
and blocks satisfy
\[
 (A,C,V)_{m+6}=(A,C,V)_m,\qquad z_{m+6}=z_m.
\]
\end{proposition}

\begin{proof}
The translation \(t\mapsto t+54\) carries \(E_m\) bijectively to
\(E_{m+6}\) and preserves each counted incidence and its multiplicity. The
equality is transmitted to the sum and to decimal reduction.
\end{proof}

Therefore, a presentation with \(12\) distinct blocks requires the reader to
preserve some information that does not factor through that observable of
period \(54\): incidence selection, effective orientation, endpoints,
multiplicity, or memory. The proposition delimits an inadmissible reduction
of the state; it does not assert that the complete dynamics has period \(54\),
nor does it by itself determine its terminal reader.
